%==============================================================================
% ASHEN — working draft (v0.1). IEEE conference style.
% Compile: pdflatex ashen.tex  (needs the IEEEtran class; available on Overleaf
% and most TeX distributions). One experiment (E1) and a remediation pilot (6.3)
% carry live data; remaining experiments are marked [PLACEHOLDER].
% Numbers are pilot-scale (3 hosts, 4 exploit types) — trends, not significance.
%==============================================================================
\documentclass[conference]{IEEEtran}
\usepackage{booktabs}
\usepackage{array}
\usepackage{xcolor}
\usepackage{amssymb}
\usepackage{pifont}
\usepackage{hyperref}
\newcommand{\todo}[1]{\textcolor{red}{\textbf{[PLACEHOLDER: #1]}}}
\newcommand{\code}[1]{\texttt{#1}}

\begin{document}

\title{ASHEN: Action-Aware Retrieval-Augmented LLM Guidance for Authorized
Penetration Testing with Closed-Loop Remediation Validation}

\author{\IEEEauthorblockN{\todo{author name}}
\IEEEauthorblockA{\todo{affiliation, email}}}

\maketitle

\begin{abstract}
Large language models (LLMs) are increasingly proposed to assist penetration
testing, but most systems stop at recommendation and are evaluated subjectively.
We present \textbf{ASHEN}, an AI-assisted penetration-testing platform that closes
the full loop --- scan, vulnerability detection, exploit validation,
retrieval-augmented attack recommendation, remediation guidance, and
\emph{re-validation of the generated fix by re-exploitation} --- under an explicit
authorization-and-audit governance model, using a \emph{local} LLM so that no
target data leaves the operator's network. We make three contributions. First, an
evaluation harness that uses ASHEN's own exploit runner as an \emph{objective
oracle}: an exploit type is ``correct'' for a target iff running it returns a
positive verdict, enabling scoring of recommendation ranking and remediation
efficacy without human judgement. Second, using this oracle we show that
\emph{naively} injecting retrieved CVE knowledge into the recommender degrades
action selection, and that an \emph{action-aware grounding} scheme --- mapping
retrieved CVEs onto the tool's available exploit types --- restores correct
ranking while keeping the recommendation grounded. Third, we define
\emph{remediation efficacy} as the fraction of confirmed vulnerabilities that no
longer validate after the generated fix is applied. In a feasibility pilot (4
remediation trials, efficacy 0.25), the loop caught three confident-but-incorrect
LLM fixes --- including one whose command \emph{executed successfully and passed
the LLM's own validation step} yet left the host exploitable --- that a
recommendation-only, human-rated, or self-validating evaluation would have
accepted. Applying the same grounding scheme to the remediation step then recovers
the failures end-to-end (efficacy 0.25$\rightarrow$0.75), unifying our two results:
grounding corrects both recommendation and remediation. We describe the protocol
for scaling to a larger evaluation.
\end{abstract}

%------------------------------------------------------------------------------
\section{Introduction}
Penetration testing is a manual, expertise-intensive process, motivating a
substantial body of work on \emph{autonomous penetration testing}~\cite{rlreview}.
Two families of decision-making methods dominate: reinforcement-learning (RL)
agents that learn attack and path-planning policies in simulated
networks~\cite{rlreview,autoptsim}, and, more recently, large language model (LLM)
agents that exploit the models' contextual semantic understanding to drive the
workflow~\cite{pentestgpt,hacksynth}. ASHEN belongs to the LLM-based family;
existing LLM-based pentest systems, however, share three weaknesses that limit
their scientific standing:
\begin{enumerate}
\item \textbf{They stop at recommendation.} Whether a suggestion \emph{works},
and whether any proposed fix actually \emph{removes} the weakness, is unmeasured.
\item \textbf{They are evaluated subjectively}, by human raters or anecdotes, so
results are hard to reproduce and easy to contest.
\item \textbf{They send sensitive data to a third-party API} --- target IPs,
detected vulnerabilities, credentials --- a confidentiality problem for real
engagements.
\end{enumerate}

ASHEN addresses all three. It runs an \emph{on-premise} LLM (no data
exfiltration), it \emph{validates} every exploit it recommends and every fix it
proposes against the live target, and it is wrapped in an authorization-gated
governance model with dual audit logs. Because ASHEN can \emph{execute} exploits,
the same execution serves as an \emph{objective oracle} for evaluation.

\noindent\textbf{Contributions.}
\begin{itemize}
\item \textbf{C1 --- An oracle-backed evaluation harness} (\S\ref{sec:method})
that scores recommendation ranking and remediation efficacy using the exploit
runner's verdict as ground truth, with no human labelling.
\item \textbf{C2 --- Action-aware retrieval grounding}
(\S\ref{sec:design},~\S\ref{sec:e1}): dumping retrieved CVEs into the prompt
biases the model toward CVE-shaped actions the tool cannot perform; mapping
retrieved CVEs onto available exploit types fixes this. \emph{(E1, n=3 hosts:
robust across Linux, Windows, and Shellshock targets where plain RAG is not; more
targets \todo{at scale}.)}
\item \textbf{C3 --- Closed-loop remediation efficacy}
(\S\ref{sec:design},~\S\ref{sec:remediation}): an end-to-end metric that
re-exploits the target after the generated fix is applied. \emph{(Pilot, n=4,
efficacy 0.25$\rightarrow$0.75 grounded: the loop caught three incorrect fixes --- including one that executed
successfully and passed its own validation yet left the host exploitable; full
study \todo{more trials}.)}
\end{itemize}

%------------------------------------------------------------------------------
\section{Background and Related Work}
\textbf{Autonomous and automated penetration testing.} A large literature casts
penetration testing as a sequential decision-making problem and applies
reinforcement learning (RL) to learn attack and path-planning policies in
simulated networks. Liu et al.~\cite{rlreview} review RL-based decision-making for
autonomous pentesting, and unified modeling frameworks such as
AutoPT-Sim~\cite{autoptsim} standardise the simulation environments these agents
train in. Such methods excel at policy learning but operate in simulation, produce
no human-readable guidance, and address neither remediation nor data
confidentiality --- gaps that motivate ASHEN's LLM-based, on-premise design.

\textbf{LLMs for offensive security.} The closest system is
PentestGPT~\cite{pentestgpt}, which integrates GPT-3.5/GPT-4 and organises a
pentest around a \emph{Pentesting Task Tree} driven by three modules (Reasoning,
Generation, Parsing). It is evaluated on 13 HackTheBox/VulnHub targets (182
sub-tasks) and reports a 228.6\% task-completion increase over its GPT-3.5
baseline. Two properties separate it from ASHEN: it is cloud/API-backed (target
data leaves the network), and it is human-in-the-loop and \emph{offensive-only} ---
the operator executes the suggested commands and feeds results back, with no
remediation-generation or fix-validation component. A parallel line shows LLM
agents can autonomously \emph{exploit}: Fang et al.\ demonstrate autonomous
website hacking~\cite{fangweb}, exploitation of real one-day CVEs with a GPT-4
agent~\cite{fangoneday}, and teams of agents exploiting zero-day
vulnerabilities~\cite{fangzeroday}; Happe and Cito study autonomous Linux
privilege escalation~\cite{happe}. Most directly comparable, HackSynth~\cite{hacksynth}
pairs a dual-module LLM agent (Planner/Summarizer) with a CTF-based \emph{evaluation
framework} and, like ASHEN, stresses robust safeguards. These drive offence but
stop at exploitation, judge success by human or capture-the-flag task-completion
rather than an execution oracle, and do not generate or re-validate remediations. We compare capabilities against PentestGPT's published
design (Table~\ref{tab:pgpt}) rather than re-running it (its native backend is a
paid GPT-4 API; a fair empirical head-to-head is future work, \S\ref{sec:remaining}).

\begin{table}[h]
\caption{ASHEN vs PentestGPT (capabilities; PentestGPT facts from~\cite{pentestgpt}).}
\label{tab:pgpt}
\centering\footnotesize
\begin{tabular}{p{2.3cm}p{2.4cm}p{2.4cm}}
\toprule
Capability & PentestGPT & ASHEN \\
\midrule
LLM backend & GPT-3.5/4 (cloud) & \textbf{local} (llama3.2) \\
Data exposure & sent to OpenAI & \textbf{on-premise} \\
Execution & human-in-the-loop & \textbf{automated} \\
Validation & none & \textbf{execution-as-oracle} \\
Grounding & task-tree context & RAG + action-aware \\
Remediation & \textbf{none} & RAG-grounded \\
Fix re-validation & \textbf{none} & \textbf{re-exploit loop} \\
\bottomrule
\end{tabular}
\end{table}

\textbf{Retrieval-augmented generation (RAG).} RAG~\cite{rag} grounds LLM output
in a retrieved corpus to reduce hallucination. We show
(\S\ref{sec:e1}) that in a \emph{tool-using} setting, naive RAG can instead
\emph{mis-direct action selection} when the knowledge representation (CVEs) is
misaligned with the action space (exploit types).

\textbf{LLM-assisted remediation.} LLMs have been applied to vulnerability
repair, e.g.\ Pearce et al.'s zero-shot vulnerability repair~\cite{pearce}, which
generates code fixes evaluated against a test suite rather than by re-exploiting a
live target. Closest to ASHEN is PenHeal~\cite{penheal}, a two-stage framework
pairing a pentest module with a remediation module; it \emph{recommends and
scores} fixes but does \textbf{not} re-validate them by re-exploiting the patched
system --- precisely the gap ASHEN's closed loop fills. To our knowledge ASHEN is
the first to score remediation by re-exploitation.

\textbf{Position of ASHEN.} ASHEN is distinguished by the combination of (i)
local LLM, (ii) execution-as-oracle evaluation, and (iii) closed-loop remediation
validation, under (iv) an explicit authorization/audit governance model.

%------------------------------------------------------------------------------
\section{ASHEN System Design}\label{sec:design}
\subsection{Overview and roles}
ASHEN separates an \emph{Analyst} (runs scans/exploits on authorised targets,
sees only their own work) from an \emph{Admin} (whitelists target IPs, approves
target requests, reads the audit log). The assessment workflow is: \emph{Scan}
(Nmap~\cite{nmap} \code{-sV -{}-script vuln}) $\rightarrow$ \emph{Vulnerability} (an
unvalidated NSE detection) $\rightarrow$ \emph{Exploit Run} (validates or refutes
it) $\rightarrow$ \emph{Attack Recommendation} (RAG-grounded LLM guidance)
$\rightarrow$ \emph{Remediation} (LLM fix guidance) $\rightarrow$ \emph{Report}.

\subsection{Exploit types as a typed action space}
ASHEN exposes a small, typed set of \emph{exploit types}, each a uniform adapter
producing a single verdict envelope (\code{ran}, \code{vulnerable},
\code{evidence}). The current set is four: \code{ssh\_brute\_force},
\code{ftp\_brute\_force} (credential family) and \code{ms17\_010\_check},
\code{shellshock\_cgi} (check family). This typed action space is what makes both
the oracle (\S\ref{sec:method}) and action-aware grounding possible.

\subsection{Retrieval grounding (and its failure mode)}
The recommender retrieves CVE context from a vector store (ChromaDB~\cite{chroma}
over sentence-transformer embeddings~\cite{sbert}) and conditions a recommendation
on it. We
compare three grounding conditions: \emph{off} (no retrieval); \emph{plain}
(retrieved CVEs concatenated into the prompt --- the baseline); and
\emph{action-aware} (each retrieved CVE is mapped onto the exploit types whose
declared \code{service} + \code{signatures} it matches; CVEs matching no available
exploit are surfaced as informational only). The corpus stays tool-agnostic;
adding an exploit type changes only that exploit's declaration, not the corpus.

\subsection{Closed-loop remediation}
For a confirmed-vulnerable finding, ASHEN generates a five-part remediation (Root
Cause / Immediate Containment / Permanent Fix / Validation / Hardening). The fix
is applied to the target, after which the \emph{same exploit is re-run}: the
vulnerability is ``remediated'' iff it no longer validates \emph{and} the service
is still reachable.

\subsection{Governance and the local-LLM choice}
Every operator action is recorded in an Audit Log; every LLM prompt/response in a
separate AI Governance Log. The LLM runs on-premise, so target data never leaves
the network --- a deliberate confidentiality choice that distinguishes ASHEN from
API-backed tools.

%------------------------------------------------------------------------------
\section{Evaluation Methodology}\label{sec:method}
\subsection{Execution as an objective oracle}
For a target and exploit type, the exploit runner returns a boolean
\code{vulnerable} verdict, treated as \emph{ground truth}: an exploit type is
\emph{relevant} iff it returns \code{vulnerable=True}. This removes human
judgement from scoring (a) whether a recommendation ranks working exploits first,
and (b) whether a remediation removed the weakness.

\subsection{Metrics}
\textbf{Precision@1 / MRR} of the recommended exploitation order against the
relevant set. \textbf{Fabrication rate}: fraction of cited CVEs not in a curated
set of CVEs real-and-applicable to the target. \textbf{Remediation efficacy}:
fraction of confirmed-vulnerable findings that, after the generated fix is
applied, no longer validate \emph{and} leave the service up (\code{fixed} vs
\code{not\_fixed} vs \code{broke\_service}).

\subsection{Harness and reproducibility}
The harness drives ASHEN's services directly (no web/auth/DB layer). Generation
is seeded (\code{temperature=0}, \code{seed=42}). A full live run is split into
isolated phases (oracle / recommend / remediate) so memory-heavy components do not
run concurrently.

\subsection{Testbed}
\emph{(Pilot, 3 hosts.)} Three targets spanning all four exploit types with
positive cases: Metasploitable~2 (Linux; vsftpd 2.3.4 backdoor~\cite{vsftpd},
OpenSSH 4.7p1, Apache 2.2.8, Samba 3.x $\rightarrow$ ssh/ftp brute-force),
Metasploitable~3 Windows Server 2008 R2 (SMBv1 $\rightarrow$
MS17-010/EternalBlue~\cite{ms17010}), and an Ubuntu~14.04 host (Apache mod\_cgi +
unpatched bash $\rightarrow$ Shellshock~\cite{shellshock}). Further scale-up
protocol in \S\ref{sec:remaining}.

%------------------------------------------------------------------------------
\section{E1: Grounding and Recommendation Quality}\label{sec:e1}
\emph{Live, n=3 hosts (all four exploit types), pilot.} We use three hosts:
Metasploitable~2 (Linux; relevant $=\{$\code{ftp\_brute\_force},
\code{ssh\_brute\_force}$\}$), Metasploitable~3 Windows Server 2008 R2
(relevant $=\{$\code{ms17\_010\_check}$\}$, EternalBlue), and an Ubuntu~14.04
Shellshock host (relevant $=\{$\code{shellshock\_cgi}$\}$). Ground truth is set per
target by the oracle; fabrication is scored against a \emph{per-target} set of
applicable CVEs, so EternalBlue counts as valid on Windows but as a fabrication on
Linux.

\begin{table}[t]
\caption{Grounding conditions across three hosts / all four exploit types (live,
seeded $\tau{=}0$, n=3). Aggregate P@1: off $=1.00$, plain $=0.67$,
action-aware $=1.00$.}
\label{tab:e1}
\centering
\begin{tabular}{llccc}
\toprule
Target & Grounding & Top-1 & P@1 & Fabr. \\
\midrule
MS2 (Lin) & off          & \code{ftp\_brute\_force} \checkmark & 1.00 & --- \\
MS2 (Lin) & plain        & \code{ms17\_010\_check} \ding{55}   & 0.00 & 0.60 \\
MS2 (Lin) & action-aware & \code{ssh\_brute\_force} \checkmark & 1.00 & 1.00$^{*}$ \\
\midrule
MS3 (Win) & off          & \code{ms17\_010\_check} \checkmark  & 1.00 & --- \\
MS3 (Win) & plain        & \code{ms17\_010\_check} \checkmark  & 1.00 & --- \\
MS3 (Win) & action-aware & \code{ms17\_010\_check} \checkmark  & 1.00 & --- \\
\midrule
Shellshock & off          & \code{shellshock\_cgi} \checkmark & 1.00 & --- \\
Shellshock & plain        & \code{shellshock\_cgi} \checkmark & 1.00 & 0.60 \\
Shellshock & action-aware & \code{shellshock\_cgi} \checkmark & 1.00 & 0.50 \\
\bottomrule
\end{tabular}
\end{table}

\noindent\textbf{Findings.}
\begin{enumerate}
\item \textbf{Plain RAG is unreliable; its failure is target-dependent.} Plain
retrieval biases the model toward SMB/CVE-heavy exploits regardless of
applicability --- \emph{wrong} on the Linux host (ranks the dead-end
\code{ms17\_010} first, P@1 $=0$) but \emph{coincidentally right} on the Windows
host (where EternalBlue truly is the answer, P@1 $=1$; Shellshock is correctly
ranked too). Aggregated it scores only P@1 $=0.67$ --- still the worst condition.
\item \textbf{Action-aware grounding is robust} (P@1 $=1.0$ across all three hosts).
Mapping retrieved CVEs onto available exploit types and quarantining unmatched
ones keeps the recommendation correct \emph{and} grounded on every host --- the
only condition achieving both; plain sacrifices ranking, \emph{off} sacrifices
grounding (it cites nothing).
\item \textbf{The per-target CVE distinction validates the design.} The same CVE
(\code{CVE-2017-0144}) is a fabrication on Linux and valid evidence on Windows,
because correctness is judged against each host's actual exposure.
\item \textbf{A residual fabrication artifact ($^{*}$, MS2).} Action-aware still
cited one off-platform CVE for a low-priority exploit, motivating
platform-applicability filtering (\S\ref{sec:discussion}).
\end{enumerate}

\noindent\textbf{Interpretation.} The mechanism is a \emph{knowledge--action-space
misalignment}: grounding a tool-using LLM on a CVE corpus biases it toward
CVE-shaped actions, including ones outside the tool's capability or inapplicable
to the host. Aligning the knowledge to the action space, and judging it
per-target, resolves this. A 5-seed $\times$ 2-temperature variance study
(\S\ref{sec:remaining}) confirms this is not a single-seed artifact: at a
task-appropriate temperature action-aware is perfectly stable (1.00 $\pm$ 0.00)
and plain RAG is the worst at every temperature. \emph{Caveat: n=3 hosts.}

%------------------------------------------------------------------------------
\section{Headline: Remediation Efficacy}\label{sec:remediation}
\emph{Live, n=4 trials across 3 hosts, pilot.} Four confirmed findings were
remediated end-to-end: ASHEN generated its guidance, \emph{its prescribed fix was
applied verbatim} to the live host (snapshotted first), and the same exploit was
re-run.

\begin{table}[t]
\caption{Remediation efficacy across three hosts (live, n=4 pilot).
\textbf{Efficacy $=0.25$ (1/4 fixed)}.}
\label{tab:remediation}
\centering
\begin{tabular}{p{1.9cm}p{2.5cm}p{1.4cm}c}
\toprule
Host / finding & ASHEN's prescribed fix & After & Outcome \\
\midrule
MS2 \code{ssh\_brute} & rotate weak creds & no creds & \textbf{fixed} \\
MS2 \code{ftp\_brute} & rotate creds (+ bogus \code{allow\_root}) & \code{ftp:<any>} logs in & \textbf{not\_fixed} \\
MS3 \code{ms17\_010} & \code{reg add ... DisableSMBv1=1} (succeeds, key present) & still vulnerable after reboot & \textbf{not\_fixed} \\
Shellshock & update Apache mod\_cgi (wrong component) & still vulnerable & \textbf{not\_fixed} \\
\bottomrule
\end{tabular}
\end{table}

\noindent\textbf{Findings.}
\begin{enumerate}
\item \textbf{The loop caught a plausible-but-wrong fix (MS2 FTP).} ASHEN asserted
weak credentials; re-exploitation showed the real cause is \emph{anonymous login}
(the \code{ftp} account accepts any password), which credential rotation cannot
address.
\item \textbf{It caught a fix that \emph{executed successfully but did nothing}
(MS3 MS17-010).} ASHEN diagnosed the root cause correctly (SMBv1) but prescribed a
non-existent registry value (\code{DisableSMBv1}; the effective control is
\code{SMB1=0}). The command returns ``completed successfully,'' the key is
verifiably present, \emph{and ASHEN's own suggested validation (a \code{reg query})
passes} --- yet after reboot the host is still EternalBlue-vulnerable. Every check
short of re-attacking reports success; only re-exploitation reveals the fix is
hollow.
\item \textbf{All three ungrounded failures misdiagnosed the fix} ---
\code{allow\_root} (non-existent), \code{DisableSMBv1} (wrong registry value),
and ``update Apache'' for Shellshock (wrong component: it is a \emph{bash} bug) ---
linking to \S\ref{sec:e1}.
\item \textbf{Where diagnosis \emph{and} command were correct, the fix worked
(MS2 SSH)} --- the positive control.
\end{enumerate}

\noindent\textbf{Interpretation.} The headline metric exposes what
recommendation-only, human-rated, \emph{or even self-validating} LLM remediation
would miss: a fix's confidence, structure, successful execution, and passing its
own validation step are all uncorrelated with whether the vulnerability is gone.
Only execution-based re-validation reveals the difference.
\emph{Caveat: n=4 trials across three hosts; not a rate estimate.}

\noindent\textbf{Grounding the remediation closes the gap (re-validated
end-to-end).} The remediation path, unlike the recommender (\S\ref{sec:e1}), was
originally \emph{ungrounded} free-form generation --- the source of the
hallucinated fixes. We applied the same grounding principle (retrieve curated,
verified fix references; constrain the LLM to them) and re-ran all four
remediations end-to-end (generate $\rightarrow$ apply verbatim $\rightarrow$
re-exploit).

\begin{table}[t]
\caption{Ungrounded vs grounded remediation, re-validated end-to-end.
Efficacy $0.25 \rightarrow 0.75$.}
\label{tab:grounded}
\centering
\begin{tabular}{lll}
\toprule
Finding & Ungrounded & Grounded \\
\midrule
MS2 \code{ssh\_brute}  & rotate creds: fixed & rotate creds: fixed \\
MS2 \code{ftp\_brute}  & \code{allow\_root}: not\_fixed & \code{anon..=NO}: not\_fixed \\
MS3 \code{ms17\_010}   & \code{DisableSMBv1}: not\_fixed & \code{SMB1=0}: \textbf{fixed} \\
Shellshock & update Apache: not\_fixed & disable CGI: \textbf{fixed} \\
\bottomrule
\end{tabular}
\end{table}

Grounding \emph{corrected the prescription on all four findings} (SSH creds,
\code{SMB1=0}, \code{anonymous\_enable=NO} for FTP replacing the invented
\code{allow\_root}, and disabling mod\_cgi for Shellshock); three of four then
re-validate as fixed. The remaining FTP
failure is now isolated to a \emph{deployment-execution} problem --- MS2's
inetd-launched vsftpd does not honour \code{anonymous\_enable=NO} on restart ---
and is no longer a prescription error: the loop confirms the fix is correct yet
ineffective on that host, which is itself the argument for execution-based
re-validation. This unifies the paper's two results: the same grounding principle
corrects both recommendation and remediation.

\noindent\textbf{Complete exploit-type coverage.} A Shellshock target (Ubuntu
14.04, Apache mod\_cgi, unpatched bash) was added so the fourth exploit type ---
previously only ever a true-negative --- is exercised end-to-end. The oracle
confirmed Shellshock RCE; the grounded remediation correctly prescribed disabling
mod\_cgi (and patching bash); applying it (\code{a2dismod cgi cgid}) re-validated
as \textbf{fixed} (injection no longer executes, web service up). Across all four
exploit types, grounded remediation efficacy is \textbf{0.75 (3/4)} ---
\code{ssh\_brute\_force}, \code{ms17\_010\_check}, \code{shellshock\_cgi} fixed;
\code{ftp\_brute\_force} not\_fixed (the deployment-execution case).

%------------------------------------------------------------------------------
\section{Variance and Remaining Evaluation}\label{sec:remaining}

\noindent\textbf{Seed/temperature variance (live).} To check the deterministic
E1 result (\S\ref{sec:e1}) is not a single-seed artifact, we re-ran all
conditions across 5 seeds and 3 targets at two temperatures. Exploit selection is
a \emph{decision} task, so a low temperature ($\tau{=}0.3$) is appropriate; we
also stress-test at $\tau{=}0.7$.

\begin{table}[h]
\caption{Precision@1 variance (mean $\pm$ stdev, 5 seeds $\times$ 3 targets).}
\centering
\begin{tabular}{lcc}
\toprule
Condition & P@1 @ $\tau{=}0.3$ & P@1 @ $\tau{=}0.7$ \\
\midrule
off          & \textbf{1.00 $\pm$ 0.00} & 0.93 $\pm$ 0.26 \\
plain        & 0.67 $\pm$ 0.47 & 0.50 $\pm$ 0.50 \\
action-aware & \textbf{1.00 $\pm$ 0.00} & 0.79 $\pm$ 0.41 \\
\bottomrule
\end{tabular}
\end{table}

\noindent Plain RAG is the worst at \emph{both} temperatures (a robust negative
result). At the task-appropriate temperature action-aware is perfectly stable
(1.00 $\pm$ 0.00), matching the ungrounded baseline while keeping its CVE
justification; its $\tau{=}0.7$ dip is a diagnosed high-temperature artifact (the
model occasionally brute-forcing open ports against the grounding), still well
above plain RAG. Fabrication is high/noisy for both RAG conditions
($\sim$0.6--0.7) and does not discriminate them, so we de-emphasise it.

\medskip
\noindent\textbf{Remaining.}
\todo{E2 --- ranking across many targets with confidence intervals.}
\todo{Headline at scale: remediation efficacy across $N\approx20$--$30$ findings,
on a testbed adding Metasploitable~3, VulnHub images, a Windows MS17-010 host, and
a Shellshock CGI host to exercise all four exploit types with positive cases.}
\todo{Model-strength comparison: weak (tinyllama) vs default (llama3.2) vs strong
(GPT-4 via API) to separate ``small model is weak'' from ``grounding scheme''.}
\textit{PentestGPT: a capability comparison is given in \S2
(Table~\ref{tab:pgpt}); a same-target empirical head-to-head (real
PentestGPT-GPT-4) is future work --- its paid GPT-4 backend and GPT-4-tuned parser
make a fair local run impractical.}

%------------------------------------------------------------------------------
\section{Discussion}\label{sec:discussion}
The headline conceptual result is the \emph{knowledge--action-space misalignment}
in tool-using RAG (\S\ref{sec:e1}): grounding helps factual accuracy but can hurt
\emph{action selection} when retrieved knowledge references capabilities the tool
lacks. Action-aware grounding is a general remedy. A residual refinement is
\emph{platform-applicability filtering}: do not attach a CVE as evidence for an
exploit when the CVE's product/platform contradicts the target fingerprint.

%------------------------------------------------------------------------------
\section{Limitations}
\emph{Few targets} (n=3 hosts for E1, n=4 trials for remediation): pilot numbers are not yet significant;
\S\ref{sec:remaining} defines scale-up. \emph{Four exploit types}: shallow
ranking depth; power must come from many targets. \emph{Manual remediation}: ASHEN
generates guidance; applying the fix is a human (or operator-authorised) step.
\emph{Weak default model}: llama3.2-3B is small.

%------------------------------------------------------------------------------
\section{Ethics and Responsible Use}
ASHEN operates only against IPs an Admin has explicitly authorised; Analysts
cannot act outside scope. All actions and LLM interactions are audit-logged. The
local-LLM design keeps target data on-premise. All experiments were conducted
against a self-hosted, intentionally-vulnerable VM owned by the authors on an
isolated network.

%------------------------------------------------------------------------------
\section{Conclusion}
ASHEN closes the penetration-testing loop from detection through validated
remediation, using execution as an objective oracle and a local LLM for
confidentiality. A three-host pilot demonstrates that action-aware retrieval
grounding corrects a failure mode of naive RAG in tool-using agents, and that
closed-loop re-validation catches confident-but-incorrect remediations. The
multi-target scale-up is in progress.

%------------------------------------------------------------------------------
\begin{thebibliography}{9}
\bibitem{pentestgpt} G.~Deng \emph{et al.}, ``PentestGPT: An LLM-empowered
Automatic Penetration Testing Tool,'' \emph{USENIX Security}, 2024.
arXiv:2308.06782.
\bibitem{rag} P.~Lewis \emph{et al.}, ``Retrieval-Augmented Generation for
Knowledge-Intensive NLP Tasks,'' \emph{NeurIPS}, 2020.
\bibitem{fangweb} R.~Fang \emph{et al.}, ``LLM Agents can Autonomously Hack
Websites,'' 2024. arXiv:2402.06664.
\bibitem{fangoneday} R.~Fang \emph{et al.}, ``LLM Agents can Autonomously Exploit
One-day Vulnerabilities,'' 2024. arXiv:2404.08144.
\bibitem{fangzeroday} R.~Fang \emph{et al.}, ``Teams of LLM Agents can Exploit
Zero-Day Vulnerabilities,'' 2024. arXiv:2406.01637.
\bibitem{happe} A.~Happe and J.~Cito, ``LLMs as Hackers: Autonomous Linux
Privilege Escalation Attacks,'' 2023. arXiv:2310.11409.
\bibitem{rlreview} J.~Liu, Y.~Zhang, S.~Zhou, J.~Yang, Y.~Lu, and X.~Zhong,
``Autonomous Penetration Testing using Reinforcement Learning: A Review and
Perspectives,'' \emph{Expert Systems with Applications}, 2025.
\bibitem{autoptsim} Y.~Wang \emph{et al.}, ``A Unified Modeling Framework for
Automated Penetration Testing (AutoPT-Sim),'' 2025. arXiv:2502.11588.
\bibitem{hacksynth} L.~Muzsai, D.~Imolai, and A.~Luk\'acs, ``HackSynth: LLM Agent
and Evaluation Framework for Autonomous Penetration Testing,'' 2024. arXiv:2412.01778.
\bibitem{penheal} J.~Huang and Q.~Zhu, ``PenHeal: A Two-Stage LLM Framework for
Automated Pentesting and Optimal Remediation,'' \emph{Workshop on Autonomous
Cybersecurity}, 2024. arXiv:2407.17788.
\bibitem{pearce} H.~Pearce \emph{et al.}, ``Examining Zero-Shot Vulnerability
Repair with Large Language Models,'' \emph{IEEE S\&P}, 2023.
\bibitem{nmap} G.~F.~Lyon, \emph{Nmap Network Scanning: The Official Nmap Project
Guide to Network Discovery and Security Scanning}. Insecure.Com LLC, 2009.
\bibitem{sbert} N.~Reimers and I.~Gurevych, ``Sentence-BERT: Sentence Embeddings
using Siamese BERT-Networks,'' \emph{EMNLP}, 2019.
\bibitem{chroma} Chroma, ``Chroma: the open-source embedding database,''
\url{https://www.trychroma.com}, 2023.
\bibitem{ms17010} Microsoft, ``Security Bulletin MS17-010: Security Update for
Microsoft Windows SMB Server,'' 2017. CVE-2017-0143/0144/0145.
\bibitem{shellshock} NVD, ``CVE-2014-6271 (Shellshock): GNU Bash environment
variable command injection,'' 2014.
\bibitem{vsftpd} NVD, ``CVE-2011-2523: vsftpd 2.3.4 backdoor command execution,''
2011.
\end{thebibliography}

\end{document}
